\documentclass[11pt,a4paper]{article}
\usepackage[margin=1in]{geometry}
\usepackage{amsmath,amssymb,amsthm}
\usepackage{algorithm}
\usepackage{algpseudocode}
\usepackage{booktabs}
\usepackage{hyperref}

\theoremstyle{definition}
\newtheorem{definition}{Definition}[section]
\newtheorem{theorem}{Theorem}[section]
\newtheorem{lemma}[theorem]{Lemma}
\newtheorem{remark}{Remark}[section]

\title{\textbf{Rigorous Mathematical Specification, Rectified Linear Unit Half-Space Dynamics, and Derivation of He (Kaiming) Initialization}}
\author{Team Scaibu \\ \small Email: \texttt{jaydeep.9664920749@gmail.com} \\ \small Architecture and Core Engine Team}
\date{October 2026}

\begin{document}
\maketitle

\begin{abstract}
This document presents the complete theoretical derivation and algorithmic specification of He (Kaiming) Initialization. Standard Glorot initialization fails in deep networks employing non-linear Rectified Linear Units (ReLU) because the zeroing of negative activations halves the second moment ($\mathbb{E}[\operatorname{ReLU}(z)^2] = \frac{1}{2} \operatorname{Var}(z)$), causing activation and gradient variances to collapse exponentially as $(1/2)^L$. He initialization introduces a gain factor $\gamma = \sqrt{2}$ (or $\sqrt{2 / (1 + a^2)}$ for Leaky ReLU) to achieve strict variance preservation across arbitrary depth. We formalize the half-space integration proof, formulate complete algorithmic pseudo-code, derive exact complexities, work through a numerical example, and discuss numerical bounds.
\end{abstract}

\section{Notation and Mathematical Preliminaries}

Let layer $l$ compute pre-activations $z = W x$ and post-activations $y = \phi(z)$, where $W \in \mathbb{R}^{n_{\text{out}} \times n_{\text{in}}}$ and $\phi(z) = \max(a z, z)$ with negative slope $a \in [0, 1)$.

\begin{itemize}
    \item $n_{\text{in}} \in \mathbb{N}_{\ge 1}$: Input connection fan-in.
    \item $n_{\text{out}} \in \mathbb{N}_{\ge 1}$: Output connection fan-out.
    \item $\text{mode} \in \{\text{fan\_in}, \text{fan\_out}\}$: Optimization target (forward signal preservation vs. backward gradient flow).
    \item $a \ge 0$: Negative slope coefficient for Leaky ReLU ($a = 0.0$ for standard ReLU).
    \item $\gamma = \sqrt{\frac{2}{1 + a^2}}$: Theoretical rectification gain multiplier ($\sqrt{2} \approx 1.4142$ for standard ReLU).
    \item $\sigma = \frac{\gamma}{\sqrt{\text{fan}}}$: Standard deviation for Gaussian initialization $W_{ij} \sim \mathcal{N}(0, \sigma^2)$.
    \item $\text{bound} = \sqrt{3} \sigma$: Uniform sampling half-width bound $W_{ij} \sim \mathcal{U}(-\text{bound}, \text{bound})$.
\end{itemize}

\begin{table}[h]
\centering
\caption{Summary of Mathematical Symbols for He / Kaiming Initialization}
\begin{tabular}{lll}
\toprule
\textbf{Symbol} & \textbf{Type / Domain} & \textbf{Description} \\
\midrule
$n_{\text{in}}$ & $\mathbb{N}_{\ge 1}$ & Input fan dimension \\
$n_{\text{out}}$ & $\mathbb{N}_{\ge 1}$ & Output fan dimension \\
$a$ & $\mathbb{R}_{\ge 0}$ & Negative slope parameter ($0.0$ for ReLU) \\
$\gamma$ & $\mathbb{R}_{> 0}$ & Rectification gain factor $\sqrt{2 / (1 + a^2)}$ \\
$\sigma$ & $\mathbb{R}_{> 0}$ & Normal standard deviation $\gamma / \sqrt{\text{fan}}$ \\
$\text{bound}$ & $\mathbb{R}_{> 0}$ & Uniform bound $\sqrt{3} \sigma$ \\
\bottomrule
\end{tabular}
\end{table}

\section{Problem Definition and Core Concepts}

\subsection{Intuition and Motivation: The Half-Space Collapse}
Let $z = \sum_{j=1}^{n_{\text{in}}} W_{ij} x_j$. If $W_{ij}$ and $x_j$ are independent symmetric distributions with zero mean, $z$ is symmetric around zero with $\mathbb{E}[z] = 0$. Applying the ReLU activation $\phi(z) = \max(0, z)$ zeroes out the entire negative half-space $\{z < 0\}$, which has probability mass $\frac{1}{2}$. Consequently, the second moment of the post-activation is:
\begin{equation}
\mathbb{E}[y^2] = \int_0^\infty z^2 p(z) \, dz = \frac{1}{2} \int_{-\infty}^\infty z^2 p(z) \, dz = \frac{1}{2} \operatorname{Var}(z)
\end{equation}
If weights are initialized with standard Xavier variance $\operatorname{Var}(W) = 1 / n_{\text{in}}$, then $\operatorname{Var}(z) = n_{\text{in}} \frac{1}{n_{\text{in}}} \mathbb{E}[x^2] = \mathbb{E}[x^2]$, so $\mathbb{E}[y^2] = \frac{1}{2} \mathbb{E}[x^2]$. Over an $L$-layer network, the signal decays as $(1/2)^L$ (e.g., $2^{-50} \approx 10^{-15}$ in ResNet-50). Kaiming initialization compensates by setting $\operatorname{Var}(W) = \frac{2}{n_{\text{in}}}$.

\subsection{Formal Mathematical Formulation}
\begin{definition}[He / Kaiming Parameterization]
Given fan dimension $\text{fan} \in \{n_{\text{in}}, n_{\text{out}}\}$ and negative slope $a \ge 0$:
\begin{align}
\gamma &= \sqrt{\frac{2}{1 + a^2}} \label{eq:kaiming_gain} \\
\sigma &= \frac{\gamma}{\sqrt{\text{fan}}} = \sqrt{\frac{2}{(1 + a^2) \cdot \text{fan}}} \label{eq:kaiming_sigma} \\
\text{bound} &= \sqrt{3} \cdot \sigma = \sqrt{\frac{6}{(1 + a^2) \cdot \text{fan}}} \label{eq:kaiming_bound}
\end{align}
\end{definition}

\section{Algorithmic Specification}

The computational pipeline implemented in \texttt{NnAlgoHeKaimingInit} is detailed below.

\begin{algorithm}[H]
\caption{He (Kaiming) Initialization Computation}
\begin{algorithmic}[1]
\Require Fan-in count $n_{\text{in}} \ge 1$, mode $\in \{\text{fan\_in}, \text{fan\_out}\}$, fan-out count $n_{\text{out}} \ge 1$
\Require Nonlinearity $\in \{\text{relu}, \text{leaky\_relu}\}$, negative slope $a \ge 0$
\Ensure Standard deviation $\sigma$, uniform bound $\text{bound}$, gain multiplier $\gamma$
\If{$n_{\text{in}} < 1$ \textbf{or} $n_{\text{out}} < 1$}
    \State \textbf{throw} PreconditionFailureException(``Fan dimensions must be positive'')
\EndIf
\If{$\text{nonlinearity} = \text{relu}$}
    \State $\gamma \gets \sqrt{2.0}$
\ElsIf{$\text{nonlinearity} = \text{leaky\_relu}$}
    \State $\gamma \gets \sqrt{2.0 / (1.0 + a^2)}$
\Else
    \State \textbf{throw} PreconditionFailureException(``Unrecognized nonlinearity'')
\EndIf
\State $\text{fan} \gets \begin{cases} n_{\text{in}} & \text{if } \text{mode} = \text{fan\_in} \\ n_{\text{out}} & \text{if } \text{mode} = \text{fan\_out} \end{cases}$
\State $\sigma \gets \gamma / \sqrt{\text{float}(\text{fan})}$
\State $\text{bound} \gets \sqrt{3.0} \cdot \sigma$
\State \Return $\{\text{std\_dev}: \sigma, \text{uniform\_bound}: \text{bound}, \text{gain}: \gamma\}$
\end{algorithmic}
\end{algorithm}

\section{Theoretical Guarantees and Invariance Theorem}

\begin{theorem}[Forward and Backward Variance Preservation for Leaky ReLU]
Let $z_l = W_l y_{l-1}$ and $y_l = \max(a z_l, z_l)$. If $W_l \sim \mathcal{N}\left(0, \frac{2}{(1 + a^2) n_{\text{in}}^{(l)}}\right)$, then for all layers $l \in \{1, \dots, L\}$:
\begin{equation}
\operatorname{Var}(y_l) = \operatorname{Var}(y_{l-1}) = \dots = \operatorname{Var}(x)
\end{equation}
Furthermore, setting $\operatorname{Var}(W_l) = \frac{2}{(1 + a^2) n_{\text{out}}^{(l)}}$ preserves backward gradient variance:
\begin{equation}
\operatorname{Var}\left(\frac{\partial \mathcal{L}}{\partial y_{l-1}}\right) = \operatorname{Var}\left(\frac{\partial \mathcal{L}}{\partial y_l}\right)
\end{equation}
\end{theorem}

\begin{proof}
Let $z_l \sim \mathcal{N}(0, \sigma_z^2)$. The second moment of $y_l = \phi(z_l)$ is:
\begin{align}
\mathbb{E}[y_l^2] &= \int_{-\infty}^0 a^2 z^2 p(z) \, dz + \int_0^\infty z^2 p(z) \, dz = \frac{a^2}{2} \sigma_z^2 + \frac{1}{2} \sigma_z^2 = \frac{1 + a^2}{2} \sigma_z^2
\end{align}
Since $\sigma_z^2 = n_{\text{in}} \operatorname{Var}(W) \mathbb{E}[y_{l-1}^2]$, substituting $\operatorname{Var}(W) = \frac{2}{(1+a^2) n_{\text{in}}}$ yields:
\begin{equation}
\mathbb{E}[y_l^2] = \frac{1 + a^2}{2} \left( n_{\text{in}} \frac{2}{(1 + a^2) n_{\text{in}}} \mathbb{E}[y_{l-1}^2] \right) = \mathbb{E}[y_{l-1}^2]
\end{equation}
which preserves signal energy across arbitrary network depth.
\end{proof}

\subsection{Limitations}
\begin{itemize}
    \item \textbf{Residual Stream Accumulation}: In ResNets with identity skips $x_{l+1} = x_l + F(x_l)$, variance grows linearly with depth $L$ unless residual branches are scaled by $1 / \sqrt{2}$ (see ALGO-NN-49).
\end{itemize}

\section{Complexity Analysis}

\subsection{Time and Space Complexity}
\begin{itemize}
    \item \textbf{Time}: $\mathcal{O}(1)$ for scalar statistic computation.
    \item \textbf{Space}: $\mathcal{O}(1)$ memory.
\end{itemize}

\section{Worked Numerical Example}

Consider a convolutional layer with $C_{\text{in}} = 64, K_h = 3, K_w = 3$, so $n_{\text{in}} = 64 \times 3 \times 3 = 576$, with $\text{mode} = \text{fan\_in}$, $\text{nonlinearity} = \text{relu}$.

\begin{enumerate}
    \item $\gamma = \sqrt{2.0} \approx 1.41421356$.
    \item Standard deviation:
    \begin{equation}
    \sigma = \frac{\sqrt{2}}{\sqrt{576}} = \frac{1.41421356}{24.0} \approx 0.05892557
    \end{equation}
    \item Uniform half-width bound:
    \begin{equation}
    \text{bound} = \sqrt{3} \times \sigma = 1.7320508 \times 0.05892557 \approx 0.10206207
    \end{equation}
\end{enumerate}

\section{Numerical Considerations and Edge Cases}

\begin{itemize}
    \item \textbf{Uniform Variance Identity}: Sampling $W_{ij} \sim \mathcal{U}(-\text{bound}, \text{bound})$ with $\text{bound} = \sqrt{3} \sigma$ yields $\operatorname{Var}(W) = \frac{(2 \cdot \text{bound})^2}{12} = \sigma^2$.
    \item \textbf{Mode Choice}: \texttt{fan\_in} is preferred for standard classification feedforward nets, while \texttt{fan\_out} is preferred for encoder-decoder architectures with deep transposed convolutions.
\end{itemize}

\begin{thebibliography}{9}
\bibitem{he2015kaiming} K.~He, X.~Zhang, S.~Ren, and J.~Sun, ``Delving deep into rectifiers: Surpassing human-level performance on imagenet classification,'' in \emph{IEEE International Conference on Computer Vision (ICCV)}, pp.~1026--1034, 2015.
\bibitem{glorot2010xavier} X.~Glorot and Y.~Bengio, ``Understanding the difficulty of training deep feedforward neural networks,'' in \emph{AISTATS}, 2010.
\end{thebibliography}

\end{document}
